\documentclass{beamer}
%% fink install texlive
\usetheme{Warsaw}
\usepackage{amsmath,amssymb}
\usepackage{stmaryrd}
\usepackage{graphicx}
%\usepackage{haskell}
\usepackage{code}
%\usepackage{proof}
\usepackage{theorem}
\usepackage{pstricks} 


\usepackage{tikz}
\usetikzlibrary{shapes.geometric, arrows, shapes}

\tikzstyle{startstop} = [rectangle, rounded corners, minimum width=1.5cm, minimum height=0.5cm,text centered, draw=black]
\tikzstyle{io} = [trapezium, trapezium left angle=70, trapezium right angle=110, minimum width=1.5cm, minimum height=0.5cm, text centered, draw=black]
\tikzstyle{process} = [rectangle, minimum width=1.5cm, minimum height=0.5cm, text centered, draw=black]
\tikzstyle{decision} = [diamond, minimum width=1.5cm, minimum height=0.5cm, text centered, draw=black, aspect=3]
\tikzstyle{arrow} = [thick,->,>=stealth]
\tikzstyle{doublearrow} = [thick,<->,>=stealth]
\tikzstyle{db} = [cylinder, shape border rotate=90, draw,minimum
height=1.5cm, minimum width=2cm, aspect=0.3]
\tikzstyle{www} = [cloud, cloud puffs=15.7, cloud ignores
aspect,minimum width=5cm, minimum height=2cm, align=center, draw]
\tikzstyle{blank} = [rectangle] 
\tikzstyle{container} = [rectangle, rounded corners, minimum
width=7cm, minimum height=7cm,text centered, draw=black]

\tikzstyle{darrow} = [thick,<->,>=stealth]

\usepackage{listings}


% Martin's macros
   

% groups

%%\newtheorem{lemma}{Lemma}
%%\newtheorem{definition}{Definition}
%%\newtheorem{theorem}{Theorem}
%%\newtheorem{corollary}{Corollary}
%%\newtheorem{remark}{Remark}
%%\newtheorem{conjecture}{Conjecture}
%%\newtheorem{example}{Example}

% arrays

\newcommand{\bi}{\begin{array}[t]{@{}l@{}}}
\newcommand{\ei}{\end{array}}
\newcommand{\ba}{\begin{array}}
\newcommand{\ea}{\end{array}}
\newcommand{\bda}{\[\ba}
\newcommand{\eda}{\ea\]}
\newcommand{\bp}{\begin{quote}\tt\begin{tabbing}}
\newcommand{\ep}{\end{tabbing}\end{quote}}
\newcommand{\set}[1]{\left\{
    \begin{array}{l}#1
    \end{array}
  \right\}}
\newcommand{\sset}[2]{\left\{~#1 \left|
      \begin{array}{l}#2\end{array}
    \right.     \right\}}


% symbols

\newcommand{\good}{good}
 
\newcommand{\err}{{\bf W}}
\newcommand{\VV}{{\cal V}} 
\newcommand{\GT}{{\cal T}_G}
\newcommand{\KK}{{\cal K}}
\newcommand{\TT}{{\cal T}}
\newcommand{\TC}{\mbox{\it TC}}
\newcommand{\UL}{\mbox{UL}}
\newcommand{\sUL}{\mbox{\scriptsize UL}}
\newcommand{\true}{\mbox{\sf true}}
\newcommand{\Int}{\mbox{\it Int}}
\newcommand{\Bool}{\mbox{\it Bool}}
\newcommand{\SF}{\mbox{$\cal S$}}
\newcommand{\tauvec}{\bar{\tau}}
\newcommand{\tvec}{\bar{t}}
\newcommand{\kvec}{\bar{k}}
\newcommand{\xvec}{\bar{x}}
\newcommand{\muvec}{\bar{\mu}}
\newcommand{\alphavec}{\bar{\alpha}}
\newcommand{\betavec}{\bar{\beta}}
\newcommand{\gammavec}{\bar{\gamma}}
\newcommand{\thetavec}{\bar{\theta}}
\newcommand{\deltavec}{\bar{\delta}}
\newcommand{\bvec}{\bar{b}}
\newcommand{\typo}{\Gamma}
\newcommand{\tenv}{\Gamma}
\newcommand{\typoinit}{\typo_0}
\newcommand{\static}{S}
\newcommand{\dynamic}{D}
\newcommand{\sstatic}{{\scriptstyle S}}
\newcommand{\sdynamic}{{\scriptstyle D}}
\newcommand{\BTA}{\mbox{BTA}}
\newcommand{\baseshape}{\bot}
\newcommand{\TCT}{P}
\newcommand{\Haskell}{H98}
\newcommand{\TCTinit}{P_o}

% figures, rules


\newcommand{\tlabel}[1]{\mbox{(#1)}}
\newcommand{\fig}[3]
        {\begin{figure*}[t]#3\
        \caption{\label{#1}#2}\ \hrulefill\ \end{figure*}}
\newcommand{\figurebox}[1]
        {\fbox{\begin{minipage}{\textwidth} #1 \end{minipage}}}
\newcommand{\boxfig}[3]
        {\begin{figure*}\figurebox{#3\caption{\label{#1}#2}}\end{figure*}}
\newcommand{\myirule}[2]{{\renewcommand{\arraystretch}{1.2}\ba{c} #1
                      \\ \hline #2 \ea}}

% relations, operators

\newcommand{\sat}{\mbox{\it sat}}
\newcommand{\entail}{\mbox{\it entail}}
\newcommand{\unambig}{\mbox{\it unambig}}
\newcommand{\complete}{\mbox{\it complete}}

\newcommand{\projexclusion}[1]{\mbox{$\bar{\exists}#1$}}
\newcommand{\funinst}{\mbox{\it inst}}
%%\renewcommand{\inst}{\, \vdash^{\scriptstyle i} \,}
%%\newcommand{\inst}{\, \vdash \,}
\newcommand{\ileq}{\preceq}
\newcommand{\ieq}{\simeq}
\newcommand{\topannot}[1]{|#1|}
\newcommand{\gendelta}{\mbox{\it gen}_{\delta}}
\newcommand{\genbeta}{\mbox{\it gen}_{\beta}}
\newcommand{\gen}{\mbox{\it gen}}
\newcommand{\normalize}{\mbox{\it normalize}}
\newcommand{\fail}{\mbox{\it fail}}
\newcommand{\fixpt}{\mbox{${\cal F}$}}
\newcommand{\testpa}{\mbox{${\cal T}$}}
\newcommand{\addpa}{\mbox{${\cal A}$}}
\newcommand{\matchpa}{\mbox{${\cal M}$}}
\newcommand{\refmatchpa}{\mbox{${\cal R}$}}
\newcommand{\wft}[1]{\mbox{\it wft}(#1)}
\newcommand{\tv}{\mbox{\it fv}}
\newcommand{\stv}{\mbox{\it semfv}}
\newcommand{\range}{\mbox{\it range}}
\newcommand{\lub}{\wedge}
\newcommand{\tyrel}[2]{(#1 \preceq #2)}
\newcommand{\eqty}[2]{#1=#2}
%%\newcommand{\ts}{\, \vdash \,}
\newcommand{\turns}{\, \vdash \,}
\newcommand{\tsinst}{\, \vdash^i \,}
\newcommand{\tsttv}{\, \vdash^{ttv} \,}
\newcommand{\tsUL}{\, \vdash_{\mbox{\scriptsize UL}} \,}
\newcommand{\trc}{\, \vdash_{\scriptstyle \id{inf}} \,}
\newcommand{\trcp}{\, \vdash_{\scriptstyle inf}' \,}
\newcommand{\genrel}{\, \vdash^{\scriptstyle gen} \,}
\newcommand{\asubty}[2]{(#1 \leq_a #2)}
\newcommand{\ssubty}[2]{(#1 \leq_s #2)}
\newcommand{\fsubty}[2]{(#1 \leq_f #2)}
\newcommand{\subty}[2]{(#1 \leq #2)}
\newcommand{\doubleb}[1]{[\![ #1 ]\!]}
\newcommand{\id}[1]{\mbox{\it #1}}


\newcommand{\mynote}[1]{$\spadesuit${\bf #1}$\clubsuit$}

% spacing

\newcommand{\sgap}{\quad}
\newcommand{\bgap}{\quad\quad}

% reserved words

\newcommand{\mathem}{\sf}
\newcommand{\IN}{\mbox{\mathem in}}
\newcommand{\LET}{\mbox{\mathem let}}
\newcommand{\MLET}{\mbox{\mathem mlet}}
\newcommand{\LETREC}{\mbox{\mathem letrec}}
\newcommand{\FIX}{\mbox{\mathem fix}}
\newcommand{\CASE}{\mbox{\mathem case}}
\newcommand{\TYCASE}{\mbox{\mathem typecase}}
\newcommand{\ELSE}{\mbox{\mathem else}}
\newcommand{\IF}{\mbox{\mathem if}}
\newcommand{\OF}{\mbox{\mathem of}}
\newcommand{\THEN}{\mbox{\mathem then}}
\newcommand{\INST}{\mbox{\mathem instance}}
\newcommand{\OVER}{\mbox{\mathem overload}}
\newcommand{\CLASS}{\mbox{\mathem class}}
\newcommand{\WHERE}{\mbox{\mathem where}}


%%% Local Variables: 
%%% mode: latex
%%% TeX-master: "bta"
%%% End: 


\newcommand{\rev}[1]{{#1}^{\mbox{\scriptsize r}}}
\newcommand{\lquo}{\backslash}
\newcommand{\rquo}{/}
\newcommand{\diff}{-}
\newcommand{\isect}{\cap}
\newcommand{\mymid}{~\redtxt{\mid}~}
\newcommand{\lconf}[2]{#1:{\tt #2}}

\newcommand{\ignore}[1]{}

\newcommand{\magtxt}[1]{{\magenta #1}}
\newcommand{\redtxt}[1]{{\red #1}}
\newcommand{\bluetxt}[1]{{\blue #1}}
\newcommand{\greytxt}[1]{{\gray #1}}
\newcommand{\greentxt}[1]{{\green #1}}
\newcommand{\mmleq}{\leq}
\newcommand{\pow}{\^{}}
\newcommand{\venv}{\Delta}
\newcommand{\mleq}{\mbox{\tt leq}}
\newcommand{\mas}{\mbox{\tt as}}
\newcommand{\subt}{<\!:}

\newcommand{\Nturns}{\, \vdash_{\mbox{\tiny lnf}} \,}

\newcommand{\deriv}[2]{#1\backslash #2}
\newcommand{\pderiv}[2]{pd_{#2}(#1)} %% {#1\backslash_p #2}
\newcommand{\pderivFail}[2]{pd^{err}_{#2}(#1)} %% {#1\backslash_p #2}
\newcommand{\mkEmp}[1]{{#1}_{\epsilon}}
\newcommand{\strip}[1]{#1\downarrow}
\newcommand{\beb}{\begin{exampleblock}}
\newcommand{\eeb}{\end{exampleblock}}
\newenvironment{ttprog}{\begin{trivlist}\item \tt
        \begin{tabbing}}{\end{tabbing}\end{trivlist}}


\newenvironment{grammar}{%
         \begin{center} \small%
         $\begin{array}{rcll}
         }{%
         \end{array}$\end{center}\ignorespaces%
         }

\newcommand{\rem}[1]{}
\newcommand{\comment}[1]{}

\newcommand{\kl}[1]{}

\bibliographystyle{plainnat}

\begin{document}
%\lstset{language=python}
\logo{
\includegraphics[height=0.8cm]{pics/nyp-logo.eps}\vspace{230pt}
\includegraphics[height=0.6cm]{pics/sit-logo.png}\vspace{230pt}
}
\title{Introduction to Scala} 
\author{
 Kenny Zhuo Ming Lu
}
\institute[Inst.]{School of Information Technology, Nanyang Polytechnic}
\date{\today} 


\frame{\titlepage} 

%\bibliographystyle{plainnat}
%-------------------------------------------------------------------
%-------------------------------------------------------------------

\begin{frame}[fragile]
\frametitle{Learning outcomes}
By the end of this lecture, you will have some basic understanding on
\begin{itemize}
 \item Object Oriented Programming 
 \item Functional Programming 
 \begin{itemize}
 \item Algebraic data type and pattern matching
 \item Higher order functions
 \item Partial functions
 \item Implicit functions
 \item Type classes
 \item Monads
\end{itemize}
 \item Data parallelism
 \item Concurrent programming with threads and actors
\end{itemize}
\end{frame}


%-------------------------------------------------------------------
%-------------------------------------------------------------------
\begin{frame}[fragile]
\frametitle{Scala is ... }

A general purpose programming language that is 
\begin{itemize}
 \item Object Oriented and Functional
 \item with distributed programming builtin
 \item with a rich type system and a huge set of libraries and eco-systems
 \item Widely used in research institutes (e.g. Spark) and industries
   (e.g. Twitter, Linkedin, Coursera Walmart, ... )
\end{itemize}
\end{frame}

%-------------------------------------------------------------------
%-------------------------------------------------------------------
\begin{frame}[fragile]
\frametitle{Hello World in a script}
Let's say we have a Scala file {\tt Script.scala}
\beb{}
\begin{code}
println("Hello world!")
\end{code}
\eeb
To execute it
\begin{code}
$ scala Script.scala
Hello world!
\end{code}
%$
\end{frame}

%-------------------------------------------------------------------
%-------------------------------------------------------------------
\begin{frame}[fragile]
\frametitle{Hello World in a console application}
Compiling Scala program into an application
\beb{}
\begin{code}
object Main extends App 
{
    println("Hello world!")
}
\end{code}
\eeb
{\tt object} is like a singleton class.  \\
To execute it
\begin{code}
$ scalac Main.scala
$ scala Main
Hello world!
\end{code}

\end{frame}

%-------------------------------------------------------------------
%-------------------------------------------------------------------
\begin{frame}[fragile]
\frametitle{Object Oriented Programming in Scala}
It is very similar to Java
\beb{}
\begin{code}
class Person(n:String,i:String) {
   private val name:String = n
   private val id:String   = i
   def getName():String = name
   def getId():String = id
}
\end{code}
\eeb
\begin{itemize}
 \item {\tt Person} is the name of the class, {\tt
     Person(n:String,i:String)} is a constructor.
 \item {\tt private} sets the accessibility scope of the members
 \item {\tt val} introduces a value, (immutable variable)
 \item {\tt def} introduces a method defintion.
\end{itemize}
\end{frame}

%-------------------------------------------------------------------
%-------------------------------------------------------------------
\begin{frame}[fragile]
\frametitle{Object Oriented Programming in Scala}

\beb{Class Inheritence via the {\tt extends} keyword}
\begin{code}
class Student(n:String, i:String, g:Double) 
                     extends Person(n,i) {
   private var gpa = g
   def getGPA() = gpa
   def setGPA(g:Double) = {
      gpa = g
   }
}
\end{code}
\eeb
\begin{itemize}
 \item {\tt var} introduces a mutable variable
 \item Most of the type annotations are optional and will be inferred
   by the type system.
\end{itemize}
\end{frame}

%-------------------------------------------------------------------
%-------------------------------------------------------------------

\begin{frame}[fragile]
\frametitle{Object Oriented Programming in Scala}
\beb{}
\begin{code}
class Staff(n:String, i:String, sal:Double) 
                 extends Person(n,i) {
   private var salary = sal
   def getSalary() = salary
   def setSalary(sal:Double) = 
   {
      salary = sal
   }
}
\end{code}
\eeb
\end{frame}


%-------------------------------------------------------------------
%-------------------------------------------------------------------

\begin{frame}[fragile]
\frametitle{Object Oriented Programming in Scala}
\beb{Traits are like Java interfaces and abstract classes}
\begin{code}
trait NightOwl {
   def stayUpLate():Unit 
}

class Student(n:String, i:String, g:Double) 
      extends Person(n,i) with NightOwl {
   private var gpa = g
   def getGPA() = gpa
   def setGPA(g:Double) = {
      gpa = g
   }
   override def stayUpLate():Unit = { 
      println("woohoo") 
   }
}
\end{code}
\eeb
\end{frame}



%-------------------------------------------------------------------
%-------------------------------------------------------------------


\begin{frame}[fragile]
\frametitle{Running our first Scala program}
Scala comes with an interpretor (AKA REPL)

\begin{code}
scala> :load OOP.scala
Loading OOP.scala...
defined class Person
defined trait NightOwl
defined class Student
defined class Staff

scala> val tom = new Student("Tom", "X1235", 4.0)
tom: Student = Student@601c1dfc

scala> val jerry = new Staff("Jerry", "T0001", 500000.0)
jerry: Staff = Staff@650fbe32

scala> tom.stayUpLate
woohoo
\end{code}
\end{frame}
%-------------------------------------------------------------------
%-------------------------------------------------------------------
\begin{frame}[fragile]
\frametitle{Functional Programming in Scala}

OOP is useful, but not Interesting, let's consider the FP features in Scala.


\end{frame}
%-------------------------------------------------------------------
%-------------------------------------------------------------------
\begin{frame}[fragile]
\frametitle{Algebraic Data Type}
Recall the Haskell example
\beb{}
\begin{code}
data Exp =  Val Int |  Plus Exp Exp 

simp :: Exp -> Exp 
simp e@(Val v) = e
simp (Plus (Val 0) e2) = e2
simp (Plus e1 e2) = Plus (simp e1) (simp e2)
\end{code}
\eeb
\end{frame}
%-------------------------------------------------------------------
%-------------------------------------------------------------------
\begin{frame}[fragile]
\frametitle{Algebraic Data Type}
Let's rewrite in GADT style
\beb{}
\begin{code}
data Exp where
    Val  :: Int -> Exp
    Plus :: Exp -> Exp -> Exp

simp :: Exp -> Exp 
simp e@(Val v) = e
simp (Plus (Val 0) e2) = e2
simp (Plus e1 e2) = Plus (simp e1) (simp e2)
\end{code}
\eeb
\end{frame}
%-------------------------------------------------------------------
%-------------------------------------------------------------------
\begin{frame}[fragile]
\frametitle{Algebraic Data Type} 
In Scala, this is how we define a data type
\beb{}
\begin{code}
sealed trait Exp
case class Val(v:Int) extends Exp
case class Plus(e1:Exp, e2:Exp) extends Exp

def simp(e:Exp):Exp = e match {
   case Val(v) => e
   case Plus(Val(0), e2) => e2
   case Plus(e1,e2) => Plus(simp(e1), simp(e2))
}
\end{code}
\eeb
\begin{itemize}
 \item {\tt sealed} requires all extensions to the trait {\tt Exp} must be
   declared in the same module.
 \item {\tt case} makes a class ``pattern-matchable''
\end{itemize}
\end{frame}

%-------------------------------------------------------------------
%-------------------------------------------------------------------
\ignore{
\begin{frame}[fragile]
\frametitle{Algebraic Data Type}
\beb{}
\begin{code}
def simpFix(e:Exp):Exp = {
   val e2 = simp(e)
   if (e == e2) {
      e
   } 
   else {
      simpFix(e2)
   }
}
\end{code}
\eeb
\end{frame}
}

%-------------------------------------------------------------------
%-------------------------------------------------------------------

\begin{frame}[fragile]
\frametitle{List}
\beb{}
\begin{code}
sealed trait List[A]
case class Nil[A]() extends List[A]
case class Cons[A](x:A, xs:List[A]) extends List[A]

val x = Cons(1,Cons(2,Nil()))

def length[A](l:List[A]):Int = l match {
   case Nil() => 0
   case Cons(_,xs) => 1 + length(xs)
}
\end{code}
\eeb
Observation
\begin{itemize}
 \item {\tt A} is a type variable (AKA generics)
 \item Parametric polymorphism
 \item {\tt Nil()} takes no argument, it's always (almost) singleton. 
\end{itemize}
\end{frame}
%-------------------------------------------------------------------
%-------------------------------------------------------------------

\begin{frame}[fragile]
\frametitle{List}

\beb{Mixing subtying with parametric polymorphism}
\begin{code}
sealed trait List[+A]
case object Nil extends List[Nothing]
case class Cons[A](x:A, xs:List[A]) extends List[A]

val x = Cons(1,Cons(2,Nil))

def length[A](l:List[A]):Int = l match {
   case Nil => 0
   case Cons(_,xs) => 1 + length(xs)
}
\end{code}
\eeb
\begin{itemize}
 \item {\tt case object} introduces a singleton value constructor.
 \item {\tt +} indicates that {\tt A} is a covariant type parameter,
   hence {\tt List[Nothing] <: List[A]} for any {\tt A} .
 \item {\tt Nothing} is the bottom type in the sub-typing hierachy.
\end{itemize}
\end{frame}

%-------------------------------------------------------------------
%-------------------------------------------------------------------

\begin{frame}[fragile]
\frametitle{List}

List is a builtin data type comes with Scala library
\beb{}
\begin{code}
val x = List(1,2)

def length[A](l:List[A]):Int = l match {
   case Nil => 0
   case (_::xs) => 1 + length(xs)
}
\end{code}
\eeb
\end{frame}

%-------------------------------------------------------------------
%-------------------------------------------------------------------

\begin{frame}[fragile]
\frametitle{List Comprehension}
\beb{}
\begin{code}
val nats:List[Int] = List(1,2,3,4)
val even_plus_1:List[Int] = for (nat <- nats
                                 if nat % 2 == 0
                                ) yield (nat+1)
\end{code}
\eeb
\begin{itemize}
 \item {\tt for} ... {\tt yield} denotes a sequence comprehension. 
 \item C.f. set builder notation in math,  
$$
\{ nat + 1 | nat \in nats \wedge nat~ {\tt mod}~ 2 = 0 \}
$$
\end{itemize}
\end{frame}

%-------------------------------------------------------------------
%-------------------------------------------------------------------

\begin{frame}[fragile]
\frametitle{For loop vs comprehension}
Note that {\tt for} ... {\tt yield} denotes a sequence comprehension. 
\beb{}
\begin{code}
val nats:List[Int] = List(1,2,3,4)
val even_plus_1:List[Int] = for (nat <- nats
                                 if nat % 2 == 0
                                ) yield (nat+1)
\end{code}
\eeb
While {\tt for} introduces a for loop.
\beb{}
\begin{code}
for (i <- 1 to 10) { println(i) }
\end{code}
\eeb
%
prints the numbers 1 to 10.
\beb{}
\begin{code}
for (nat <- nats if nat % 2 == 0) { println(nat) }
\end{code}
%
\eeb
prints the numbers 2 and 4.
\end{frame}



%-------------------------------------------------------------------
%-------------------------------------------------------------------
\begin{frame}[fragile]
\frametitle{List operations}
\beb{Length}
\begin{code}
val l1 = List(1,2,3)
val len = l1.length // 3
\end{code}
\eeb
\beb{Concatenation}
\begin{code}
val l2 = List(-1,-2,-3)
val l3 = l1 ++ l2 // List(1,2,3,-1,-2,3)
\end{code}
\eeb
\beb{Reverse}
\begin{code}
val l4 = l1.reverse // List(3,2,1)
\end{code}
\eeb
\end{frame}


%-------------------------------------------------------------------
%-------------------------------------------------------------------
\begin{frame}[fragile]
\frametitle{List operations}
\beb{min/max}
\begin{code}
// recall l1 = List(1,2,3)
val min = l1.min // 1
val max = l1.max // 3
\end{code}
\eeb
\beb{sorting}
\begin{code}
// recall l3 = List(1,2,3,-1,-2,3)
val sl3 = l3.sortWith(_<_) // List(-3, -2, -1, 1, 2, 3)
// we will look into _<_ in a short while.
\end{code}
\eeb
\beb{sub-list}
\begin{code}
val sl3a = sl3.take(3) // List(-3,-2,1)
val sl3b = sl3.drop(3) // List(1,2,3)
\end{code}
\eeb
\end{frame}

%-------------------------------------------------------------------
%-------------------------------------------------------------------
\begin{frame}[fragile]
\frametitle{Map and Fold}
{\tt l.map(f)} returns a new list by applying the unary function {\tt f} to each
element in {\tt l}
\beb{}
\begin{code}
val l = List(1,2,3,4)
val l2 = l.map(x => x * 2) // List(2,4,6,8)
\end{code}
\eeb
%
Note that {\tt x => x * 2} is a lambda-expression, which can be
shortened as {\tt \_ * 2}. The {\tt .} is optional.
\beb{}
\begin{code}
val l2 = l map (_*2) 
\end{code}
\eeb
\end{frame}

%-------------------------------------------------------------------
%-------------------------------------------------------------------
\begin{frame}[fragile]
\frametitle{Map and Fold}

{\tt l.foldLeft(b)(f)} agreggate all the
elements in {\tt l} with the binary function {\tt f} and the base value
{\tt b}.
\beb{}
\begin{code}
// recall l = List(1,2,3,4)
val l3 = l.foldLeft(0)((x,y) => x + y)
\end{code}
\eeb
In essence the above is computing 
\begin{code}
(((0 + 1) + 2) + 3) + 4
\end{code}
{\tt l.foldRight(b)(f)} agreggate all the
elements in {\tt l} with the binary function {\tt f} and the base value
{\tt b} with right associativity.
\beb{}
\begin{code}
val l4 = l.foldRight(0)((x,y) => x + y) 
// or val l4 =  l.foldRight(0)(_+_) 
\end{code}
\eeb
In essence the above is computing 
\begin{code}
0 + (1 + (2 + (3 + 4)))
\end{code}

\end{frame}

%-------------------------------------------------------------------
%-------------------------------------------------------------------
\begin{frame}[fragile]
\frametitle{Higher order function}
Let's consider the defintion of {\tt map}.

\beb{}
\begin{code}
object List { // this is like a package
    sealed trait List[+A] {
        def map[B](f:A=>B):List[B] = this match {
            case Nil => Nil
            case Cons(x,xs) => Cons(f(x),xs.map(f))
        }
    }
    case object Nil extends List[Nothing]
    case class Cons[A](x:A, xs:List[A]) extends List[A]
}
\end{code}
\eeb
Function {\tt f:A=>B} is used as a formal argument.
\end{frame}

%-------------------------------------------------------------------
%-------------------------------------------------------------------
\begin{frame}[fragile]
\frametitle{Higher order function}
Function composition
\beb{}
\begin{code}
val l1 = List(1,2,3)
val l2 = l1.map (((x:Int) => x*2) compose (y=>y+1)) 
//  List(4, 6, 8)
\end{code}
\eeb
{\tt (f compose g)(x)} is equivalent to  {\tt f(g(x))}
\beb{}
\begin{code}
val l3 = l1.map (((x:Int) => x*2) andThen (y=>y+1))  
// List(3, 5, 7)
\end{code}
\eeb
{\tt (f andThen g)(x)} is equivalent to  {\tt g(f(x))}

\end{frame}


%-------------------------------------------------------------------
%-------------------------------------------------------------------
\begin{frame}[fragile]
\frametitle{Higher order function}
In Scala, functions are instances/objects of trait {\tt Function}
\beb{The defintions of {\tt compose} and {\tt andThen}.}
\begin{code}
object Function { // this is like a package
    trait Function1[+A][-B] {
       def apply(x:A):B {...} 
       // apply() gives us
       // the syntactic sugar form like f(x)
       def compose[C](that:C=>A):C=>B = 
         (c:C) => this(that(c))
       def andThen[C](that:B=>C):A=>C = 
         (a:A) => that(this(a))
    }
}
\end{code}
\eeb

\end{frame}


%-------------------------------------------------------------------
%-------------------------------------------------------------------
\begin{frame}[fragile]
\frametitle{Functions vs methods}

In Scala there is a clear distinction between functions and methods.
\begin{itemize}
\item Methods are introduced by {\tt def} keyword, e.g. 
\begin{code}
scala> def f(x:Int) = x
f: (x: Int)Int
\end{code}
\item Functions are implementations of the 
{\tt Function}X traits. e.g. lambda expressions are functions. 
\begin{code}
scala> val g = (x:Int) => x
g: Int => Int = <function1>
\end{code}
\item Simiilarity: Both can be used as argument and results
\begin{code}
scala> List(1,2).map (f) == List(1,2).map (g)
res2: Boolean = true
\end{code}
\end{itemize}

\end{frame}
%-------------------------------------------------------------------
%-------------------------------------------------------------------
\begin{frame}[fragile]
\frametitle{Functions vs methods}

In Scala there is a clear distinction between functions and methods.
\begin{itemize}
\item Difference: methods do not possess combinators
\begin{code}
scala> f compose g
<console>:10: error: missing arguments for method f;
follow this method with `_' if you want to treat it as 
a partially applied function
              f compose g
              ^
scala> g compose f
res1: Int => Int = <function1>
\end{code}
\item To convert a method to a function.
\begin{code}
scala> f _
res1: Int => Int = <function1>
\end{code}
\end{itemize}

\end{frame}


%-------------------------------------------------------------------
%-------------------------------------------------------------------
\begin{frame}[fragile]
\frametitle{Partial functions}
There are times we need to deal with partial function, naturally, we
can use the {\tt Option} type in Scala (similar to {\tt Maybe} in
Haskell).
\beb{}
\begin{code}
def div(x:Int)(y:Int):Option[Int] =  {
    if (y == 0) None else Some(x/y)
}

val xs = List(0,1,2) map (x => div(2)(x))
\end{code}
\eeb
yields {\tt  List(None, Some(2), Some(1))}
\beb{}
\begin{code}
val ys = xs filter (x => x.nonEmpty) 
\end{code}
\eeb
yields {\tt List(Some(2), Some(1))}
\beb{}
\begin{code}
val zs = ys map ( x => x match { case Some(y) => y } ) 
\end{code}
\eeb
yields {\tt List(2,1)}
\end{frame}

%-------------------------------------------------------------------
%-------------------------------------------------------------------
\begin{frame}[fragile]
\frametitle{Partial functions}
Alternatively, we may use the Scala bultin PartialFunction 
\beb{}
\begin{code}
def div(x:Int): PartialFunction[Int,Int] = {
    case (y:Int) if y != 0 => x/y
}
val xs = List(0,1,2) collect ( div(2) ) 
\end{code}
\eeb
yields {\tt List(2, 1)}. \\ 
Note that we have to use the {\tt collect} method instead of {\tt map} 
\beb{}
\begin{code}
val ys = List(0,1,2) map ( div(2) ) 
\end{code}
\eeb
Will result in a {\tt scala.MatchError}
\end{frame}
%-------------------------------------------------------------------
%-------------------------------------------------------------------
\begin{frame}[fragile]
\frametitle{Partial functions}
Alternatively, we may use the Scala bultin PartialFunction 
\beb{}
\begin{code}
def div(x:Int): PartialFunction[Int,Int] = {
    case (y:Int) if y != 0 => x/y
}

val mkZero: PartialFunction[Int,Int] = {
  case (y:Int) if y == 0 => 0
}
val ys = List(0,1,2) map ( div(2) orElse mkZero ) 
\end{code}
\eeb
yields in a {\tt List(0, 2, 1)}
\end{frame}
%-------------------------------------------------------------------
%-------------------------------------------------------------------
\begin{frame}[fragile]
\frametitle{Implicit functions}
Implicit function and implicit parameter allow us to define a default behavior.
\beb{}
\begin{code}
implicit def logger(mesg:String):Unit = { 
    println("[Log] " + mesg)
}

def fib(x:Int)(implicit log:String=>Unit):Int = {
  x match  {
    case 0 => 0
    case 1 => 1
    case _ => {
      log("computing fib(" + x +")")
      fib(x-1) + fib(x-2)         
    }
  }
}
\end{code}
\eeb
\end{frame}
%-------------------------------------------------------------------
%-------------------------------------------------------------------
\begin{frame}[fragile]
\frametitle{Implicit functions}

\begin{code}
scala> fib(5)
[Log] computing fib(5)
[Log] computing fib(4)
[Log] computing fib(3)
[Log] computing fib(2)
[Log] computing fib(2)
[Log] computing fib(3)
[Log] computing fib(2)
res7: Int = 5
scala> fib(5)( x => Unit) 
res8: Int = 5
\end{code}
\end{frame}
%-------------------------------------------------------------------
%-------------------------------------------------------------------
\begin{frame}[fragile]
\frametitle{Implicit functions}
Scala resolves implicit parameters as follows,
\begin{enumerate}
\item Search in current scope
  \begin{enumerate}
  \item Implicits defined in current scope
  \item Implicits defined in Explicit imports
  \item Implicits defined in wildcard imports
  \end{enumerate}
\item Search for associated types in
  \begin{enumerate}
  \item Companion objects of a type
  \item Implicit scope of an argument's type 
  \item Outer objects for nested types
  \item Other dimensions
  \end{enumerate}
\end{enumerate}
\end{frame}
%-------------------------------------------------------------------
%-------------------------------------------------------------------
\begin{frame}[fragile]
\frametitle{Type classes}

Type classes was introduced in Haskell to allow mixing parametric
polymorphism and ad-hoc polymorphism.
\beb{Haskell Type Class Example}
\begin{code}
class Eq a where 
  (==) :: a -> a -> Bool

instance Eq Int where
  (==) x y = eqInt x y

instance Eq a => Eq [a] where
  (==) [] [] = True
  (==) (x:xs) (y:ys) = (x == y) && (xs == ys)
  (==) _ _ = False

check :: (Eq a) => [a] -> [a] -> Bool
check xs ys = xs == ys
\end{code}
\eeb{}

\end{frame}
%-------------------------------------------------------------------
%-------------------------------------------------------------------
\begin{frame}[fragile]
\frametitle{Type classes}
In Scala, type classes are ``encoded'' using traits and implicit
definitions, i.e. {\tt implicit val}, {\tt implicit def} and {\tt
  implicit object}
\beb{}
\begin{code}
trait Eq[A] {  def eq(x:A,y:A):Boolean }
implicit val eqInt = new Eq[Int] {
    def eq(x:Int,y:Int):Boolean =  x == y
}
implicit def eqList[A](implicit ev:Eq[A]) = 
  new Eq[List[A]] {
    def eq(xs:List[A],ys:List[A]):Boolean = 
      (xs,ys) match {
        case (Nil,Nil) => true
        case (x::xs2, y::ys2) => ev.eq(x,y) && eq(xs2,ys2)
        case (_, _) => false
     }
}
\end{code}
\eeb
\end{frame}
%-------------------------------------------------------------------
%-------------------------------------------------------------------
\begin{frame}[fragile]
\frametitle{Type classes}
By resolving implicit parameter the compiler builds the ``evidence''
{\tt ev} AKA the dictionary.
\beb{}
\begin{code}
def check[A](x:List[A], y:List[A])
   (implicit ev:Eq[List[A]]) = ev.eq(x,y) 
\end{code}
\eeb
\begin{code}
scala> check(List(1,2,3), List(1,2,3))
res6: Boolean = true
\end{code}
The above builds the evidence {\tt ev = eqList(eqInt)}

\end{frame}
%-------------------------------------------------------------------
%-------------------------------------------------------------------
\begin{frame}[fragile]
\frametitle{Type classes}
For real world application, we do not define all the type classes from
the scratch. We use a library project called {\tt scalaz}.
\url{https://github.com/scalaz/scalaz} which provides many useful type classes.
\end{frame}

%-------------------------------------------------------------------
%-------------------------------------------------------------------
\begin{frame}[fragile]
\frametitle{Monads}
A vague analogy 
\begin{itemize}
\item Algebraic data types define abstractions of run-time data (or
  values), i.e. what are the possible values.
\item Monads define abstractions of {\em computations}, i.e. what the
  possible computations.
\end{itemize}
\end{frame}

%-------------------------------------------------------------------
%-------------------------------------------------------------------
\begin{frame}[fragile]
\frametitle{Monads}
Let's model safe arithmetic computation
\beb{}
\begin{code}
sealed trait SafeVal[+T] 
case class Val[T](x:T) extends SafeVal[T]
case class Err(msg:String) extends SafeVal[Nothing]

type SafeInt = SafeVal[Int]
\end{code}
\eeb{}
\end{frame}

%-------------------------------------------------------------------
%-------------------------------------------------------------------
\begin{frame}[fragile]
\frametitle{Monads}
\beb{}
\begin{code}
def add(x:SafeInt, y:SafeInt):SafeInt = x match {
  case Err(msg) => Err(msg)
  case Val(i) => y match {
    case Err(msg) => Err(msg)
    case Val(j) => Val(i + j)
  }
}
def div(x:SafeInt, y:SafeInt):SafeInt = x match  {
  case Err(msg) => Err(msg)
  case Val(i) => y match {
    case Err(msg) => Err(msg)
    case Val(j) => { 
        if (j == 0) {Err("div by 0")} else {Val (i/j)}
    }
  }
}
\end{code}
\eeb{}
\end{frame}

%-------------------------------------------------------------------
%-------------------------------------------------------------------
\begin{frame}[fragile]
\frametitle{Monads}
We spot many similar code structure, let's rewrite them in terms of
some higher order combinators {\tt pnt} and {\tt bnd}.
\beb{}
\begin{code}
def pnt[A](x: => A):SafeVal[A] = Val(x)
def bnd[A,B](sa: SafeVal[A])
  (f: A => SafeVal[B]): SafeVal[B] = sa match {
  case Err(msg) => Err(msg)
  case Val(a)   => f(a)
}
\end{code}
\eeb{}
\begin{itemize}
\item Note that {\tt x:=> A} defines a {\em lazy} formal parameter {\tt x}
of type {\tt A}.
\item {\tt pnt} ``lifts'' / ``injects'' a value into the {\tt SafeVal}
  object.
\item {\tt bnd} ``composes'' the previous {\tt SafeVal} with the
  subsequent computation.
\end{itemize}
\end{frame}
%-------------------------------------------------------------------
%-------------------------------------------------------------------
\begin{frame}[fragile]
\frametitle{Monads}
\beb{}
\begin{code}
def add(x:SafeInt, y:SafeInt):SafeInt = {
  bnd(x)(x1 => { 
    bnd(y)(y1 => pnt(x1+y1))
  })
}
def div(x:SafeInt, y:SafeInt):SafeInt = {
  bnd(x)( x1 => {
    bnd(y)( y1 => {
      if (y1 == 0) { Err("div by 0") }
      else { pnt(x1 / y1) }
    })
  })
}
\end{code}
\eeb{}
\end{frame}

%-------------------------------------------------------------------
%-------------------------------------------------------------------
\begin{frame}[fragile]
\frametitle{Monads}
If we use type classes {\tt Monad} and {\tt MonadPlus} provided by
scalaz, we can make it even more concise. 
\beb{}
\begin{code}
implicit object safeValMonadPlus 
  extends MonadPlus[SafeVal] {
  override def point[A](a: => A):SafeVal[A] = Val(a)
  override def bind[A, B](sa: SafeVal[A])
  (f: A => SafeVal[B]): SafeVal[B] = sa match {
    case Err(msg) => Err(msg)
    case Val(a)   => f(a)
  } 
 ... 
}
\end{code}
\eeb{}
\end{frame}
%-------------------------------------------------------------------
%-------------------------------------------------------------------

\begin{frame}[fragile]
\frametitle{Monads}
\beb{}
\begin{code}
def add(x:SafeInt,y:SafeInt):SafeInt = {
  for ( x1 <- x
  ; y1 <- y
  ) yield (x1 + y1)
}
def div(x:SafeInt,y:SafeInt):SafeInt = {
  for ( x1 <- x
  ; y1 <- y
  ; if y1 != 0
  ) yield (x1 / y1)
}
\end{code}
\eeb{}
Roughly speaking, {\tt for (x <- e1 ; e2)} is desugared to 
{\tt bind(e1)(x => e2)}
\end{frame}

%-------------------------------------------------------------------
%-------------------------------------------------------------------
\begin{frame}[fragile]
\frametitle{Monads}
There many types of commonly used monads
\begin{itemize}
 \item List monad - list traversal and aggregate computation.
 \item State monad - computation with explicit state passing
 \item Reader monad - input scanning computation
\end{itemize}
Monads are composable via monad trasformer. 
\end{frame}
%-------------------------------------------------------------------
%-------------------------------------------------------------------
\begin{frame}[fragile]
\frametitle{Data Parallelism}
Scala has built-in support of data parallelism. For instance,
\beb{}
\begin{code}
object Fib extends App {
  def fib(x:Int):Int = x match {
    case 0 => 0
    case 1 => 1
    case n => fib(n-1) + fib(n-2)
  }
  println((30 to 40).toList.map(fib))
}
\end{code}
\eeb
\begin{code}
$ time scala Fib
List(832040, 1346269, 2178309, ... )
real	0m1.621s
user	0m1.624s
sys	0m0.059s
\end{code}
%$
\end{frame}

%-------------------------------------------------------------------
%-------------------------------------------------------------------
\begin{frame}[fragile]
\frametitle{Data Parallelism}
{\tt .par} converts a sequentials data collection into a paralleled
counter part. {\tt .seq} has the opposite effect.
\beb{}
\begin{code}
object ParFib extends App {
  def fib(x:Int):Int = x match {
    case 0 => 0
    case 1 => 1
    case n => fib(n-1) + fib(n-2)
  }
  println((30 to 40).toList.par.map(fib))
}
\end{code}
\eeb
\begin{code}
$ time scala ParFib
List(832040, 1346269, 2178309, ... )
real	0m0.891s
user	0m1.901s
sys	0m0.061s
\end{code}
%$
\end{frame}


%-------------------------------------------------------------------
%-------------------------------------------------------------------
\begin{frame}[fragile]
\frametitle{Concurrent Programming}

\begin{itemize}
\item Multi-threading
\item Async computation
\item Actors
\end{itemize}

\end{frame}

%-------------------------------------------------------------------
%-------------------------------------------------------------------
\begin{frame}[fragile]
\frametitle{Concurrent Programming}

Multi-threading in Scala is similar to Java
\beb{}
\begin{code}
  object ThreadFib extends App {
    def fib(x:Int):Int = x match {
      case 0 => 0
      case 1 => 1
      case n => fib(n-1) + fib(n-2)
    }

    for (i <- 30 to 40) { 
      new Thread(new Runnable {
        def run() {
          println(fib(i))
        }
      }).start
    }
  }
\end{code}
\eeb
\end{frame}


%-------------------------------------------------------------------
%-------------------------------------------------------------------
\begin{frame}[fragile]
\frametitle{Actors}
Actor model is a concurrent programming model. 
\begin{itemize}
 \item Actors are primitive objects with specific roles and functionalities
 \item Actors (objects) are coordinating with each other via message passing
\end{itemize}
\end{frame}

%-------------------------------------------------------------------
%-------------------------------------------------------------------
\begin{frame}[fragile]
\frametitle{Actors}
Let's consider a ping pong simulation using Scala actors

\begin{itemize}
\item There are two actors {\tt Ping} and {\tt Pong}. 
\item {\tt Ping} keeps an internal counter
   \begin{itemize}
      \item when it receives a {\tt start} message, 
        it increments the counter and send a {\tt ping} message to {\tt Pong}
       \item when it receives a {\tt pong} message, it increments the
         counter. 
           \begin{itemize} 
               \item if the counter $>$ 99, it sends a {\tt stop}
                 message to {\tt pong} then terminates.
               \item  otherwise, sends a {\tt ping} message to {\tt Pong}.
           \end{itemize}
   \end{itemize}
\item {\tt Pong} 
  \begin{itemize}
    \item when it receives a {\tt stop} message, it terminates.
    \item when it receives a {\tt ping} message, it sends a {\tt pong}
      message to {\tt Ping}.
  \end{itemize}
\end{itemize}

\beb{}
\begin{code}
case object PingMessage
case object PongMessage
case object StartMessage
case object StopMessage
\end{code}
\eeb
\end{frame}


%-------------------------------------------------------------------
%-------------------------------------------------------------------
\begin{frame}[fragile]
\frametitle{Actors}
\beb{}
\begin{code}
class Ping(pong: ActorRef) extends Actor {
  var count = 0
  def incrementAndPrint = {count += 1; println("ping")}
  def receive = {
    case StartMessage => 
      incrementAndPrint
      pong ! PingMessage
    case PongMessage => 
      incrementAndPrint
      if (count > 99) {
        sender ! StopMessage
        println("ping stopped")
        context.stop(self)
      } 
      else { sender ! PingMessage }
  }
}
\end{code}
\eeb
\end{frame}

%-------------------------------------------------------------------
%-------------------------------------------------------------------
\begin{frame}[fragile]
\frametitle{Actors}

\beb{}
\begin{code}
class Pong extends Actor {
  def receive = {
    case PingMessage =>
      println("  pong")
      sender ! PongMessage
    case StopMessage =>
      println("pong stopped")
      context.stop(self)
  }
}
\end{code}
\eeb
\end{frame}
%-------------------------------------------------------------------
%-------------------------------------------------------------------
\begin{frame}[fragile]
\frametitle{Actors}

\beb{}
\begin{code}

object PingPongTest extends App {
  val system = ActorSystem("PingPongSystem")
  val pong = system.actorOf(Props[Pong], 
  name = "pong")
  val ping = system.actorOf(Props(new Ping(pong)), 
  name = "ping")
  ping ! StartMessage
}
\end{code}
\eeb
\end{frame}

%-------------------------------------------------------------------
%-------------------------------------------------------------------
\begin{frame}[fragile]
\frametitle{Actors}
The ping pong example is pretty contrived, let's consider a less
contrived example in which we try to compute {\tt fib(n)} in
parallel. The idea is as follows,
\begin{itemize}
\item if n is less than 10, we compute fib(n) locally
\item if n is greater than or equal to 10, we create two fib
  actors. The first actor computes {\tt fib(n-1)} and the second actor
  computes {\tt fib(n-2)}. The two actors can run in parallel if there
  is enough threads. When both actors revert with some results, we sum up the results.
\end{itemize}
Let's define the messages for communication as follows,
\beb{}
\begin{code}
  case class FibRequest(n:Int)
  case class FibResult(n:Int)
  case object Stop
\end{code}
\eeb
\end{frame}

%-------------------------------------------------------------------
%-------------------------------------------------------------------
\begin{frame}[fragile]
\frametitle{Actors}
\beb{}
{\scriptsize
\begin{code}
class FibActor extends Actor {
  def receive = {
    case Stop => context.stop(self)
    case FibRequest(n) if n < 10 => sender ! FibResult(fib(n))
    case FibRequest(n) if n >= 10 => {
      val child1 = context.actorOf(Props[FibActor]) 
      val child2 = context.actorOf(Props[FibActor])
      val future1:Future[Any] = child1 ? FibRequest(n-1) // ``ask'' op
      val future2:Future[Any] = child2 ? FibRequest(n-2)
      val dest = sender
      val future1and2 = for ( // Future[T] is a monad
        FibResult(v1) <- future1;
        FibResult(v2) <- future2 
      ) yield v1+v2
      future1and2 onComplete {
        case Success(v) => 
        child1 ! Stop; child2 ! Stop
        dest ! FibResult(v) // can't use sender
        case Failure(e) => println("failed!" + e.toString)
      }
    }
  }
}
\end{code}
}
\eeb
\end{frame}

%-------------------------------------------------------------------
%-------------------------------------------------------------------
\begin{frame}[fragile]
\frametitle{Actors}
\beb{}
\begin{code}
def main(args:Array[String]):Unit = {
  val system = ActorSystem("FibSystem")
  val fibActor = system.actorOf(Props[FibActor])
  val n = if (args.length > 0) args(0).toInt else 10
  val f = fibActor ? FibRequest(n)
  f onComplete {
    case Success(v) => { println(v) ; fibActor ! Stop}
    case Failure(e) => println("failed!" + e.toString)
  }
}
\end{code}
\eeb
\end{frame}


%-------------------------------------------------------------------
%-------------------------------------------------------------------
\begin{frame}[fragile]
\frametitle{In Summary}
We've covered
\begin{itemize}
 \item OOP in Scala
 \item Functional Programming in Scala which includes,
  \begin{itemize}
  \item Algebraic data type and pattern matching
  \item Higher order functions
  \item Partial functions
  \item Implicit functions
  \item Type classes
  \item Monads
  \end{itemize}
 \item Parallel programming in Scala
 \item Concurrent programming in Scala with 
     \begin{itemize}
        \item threadings
        \item actors
     \end{itemize}
\end{itemize}
Latex source and example codes can be found here \\
\url{https://github.com/luzhuomi/learn_you_a_scala_for_great_good}
\end{frame}

\end{document}

